\documentclass[12pt]{article}
% Préambule commun à tous les documents extraits de « La Revue CAPES »
\usepackage[T1]{fontenc}
\usepackage[utf8]{inputenc}
\usepackage{lmodern}
\usepackage{amsmath,amssymb,amsthm,amsfonts,mathrsfs}
\usepackage{setspace}
\usepackage{listings}
\usepackage{pgfplots}
\pgfplotsset{compat=1.18}
\usepackage{longtable}
\usepackage{adjustbox}
\usepackage{lettrine}
\usepackage{tkz-tab}
\usepackage{changepage}
\usepackage{pdflscape}
\usepackage{graphicx}
\usepackage{tikz}
\usepackage{xcolor}
\usepackage[a4paper,top=2cm,bottom=2.5cm,left=2.5cm,right=2cm]{geometry}
\usepackage{fancyhdr}
\usepackage{draftwatermark}
\SetWatermarkText{La RevueCapes}
\SetWatermarkLightness{0.9}
\SetWatermarkScale{2}
\usepackage{hyperref}
\hypersetup{colorlinks=true,linkcolor=blue!50!black,urlcolor=blue!50!black}

\definecolor{revue}{RGB}{20,60,120}

\pagestyle{fancy}
\fancyhf{}
\renewcommand{\headrulewidth}{0.4pt}
\setlength{\headheight}{15pt}

% \entete{Matière}{Titre}{Type de document}
\newcommand{\entete}[3]{%
  \fancyhead[L]{\small\textsc{La Revue CAPES} -- #1}%
  \fancyhead[R]{\small #2}%
  \fancyfoot[C]{\thepage}%
  \thispagestyle{plain}%
  \begin{center}
    {\color{revue}\rule{\textwidth}{1.2pt}}\\[4pt]
    {\small\textsc{Concours directs CAP-CEG / CAPES -- Burkina Faso}}\\[6pt]
    {\LARGE\bfseries\color{revue} #2}\\[6pt]
    {\large\itshape #1 -- #3}\\[2pt]
    {\color{revue}\rule{\textwidth}{1.2pt}}
  \end{center}
  \vspace{0.5cm}%
}

% Encadré pour les remarques ajoutées dans les corrigés
\newcommand{\remarque}[1]{\par\medskip\noindent\fcolorbox{revue}{revue!6}{\parbox{\dimexpr\linewidth-2\fboxsep-2\fboxrule}{\textbf{\color{revue}Remarque.} #1}}\par\medskip}


\begin{document}
\entete{Mathématiques}{Corrigé détaillé -- Session 2018}{Correction}
\begin{spacing}{1.5}

\medskip
\textbf{Exercice 1}

On donne $A=\begin{pmatrix}
2 & 2& 2 \\ 
-1 & -1 & 1 \\ 
3 & 3 & 1
\end{pmatrix} $

1.a. Calculons les valeurs propres de $A$ et déduisons que A est diagonalisable.

Le polynôme caractéristique de $A$ est :

\begin{align*}
P_A(x)=det(A-xI_3)&=\begin{vmatrix}
2-x & 2& 2 \\ 
-1 & -1-x & 1 \\ 
3 & 3 & 1-x
\end{vmatrix}L_1\leftarrow L_1+L_2+L_3\\
&=\begin{vmatrix}
4-x & 4-x & 4-x \\ 
-1 & -1-x & 1 \\ 
3 & 3 & 1-x
\end{vmatrix}\\
&=(4-x)\begin{vmatrix}
1 & 1 & 1 \\ 
-1 & -1-x & 1 \\ 
3 & 3 & 1-x
\end{vmatrix}c_2\leftarrow c_2-c_1, c_3\leftarrow c_3-c_1\\
&=(4-x)\begin{vmatrix}
1 & 0 & 0 \\ 
-1 & -x & 2+x \\ 
3 & 0 & -2-x
\end{vmatrix}\\
&=x(4-x)(2+x)\begin{vmatrix}
-1 & 1 \\ 
 0 & -1
\end{vmatrix}\\
P_A(x) &= x(4-x)(2+x)
\end{align*}

Donc les valeurs propres de $A$ sont 0, -2 et 4

$A$ est diagonalisable car ses valeurs propres sont distinctes deux à deux. Ou autrement car le polynôme caractéristique est scindé en racines simples.

1.b. Déterminons une matrice inversible $P$ et une matrice diagonale $D$ tel que $A=PDP^{-1}$.

$D=\begin{pmatrix}
-2 & 0& 0 \\ 
0 & 0 & 0 \\ 
0 & 0 & 4
\end{pmatrix}$

Calculons les sous-espaces propres.

Soit $E_\lambda$ le sous-espace propre lié à la valeur propre $\lambda$

\[
E_\lambda=\{u=(x,y,z)\in \mathbb{R}^3/Au=\lambda u\}
\]

\medskip
$\bullet$ Pour $\lambda=0$, on a :

$\forall u\in E_0,Au=0\Rightarrow \begin{pmatrix}
2 & 2& 2 \\ 
-1 & -1 & 1 \\ 
3 & 3 & 1
\end{pmatrix}\begin{pmatrix}
x\\
y\\
z
\end{pmatrix}=\begin{pmatrix}
0\\
0\\
0
\end{pmatrix}\Rightarrow \begin{cases}
2x+2y+2z=0(1)\\
-x-y+z=0(2)\\
3x+3y+z=0(3)
\end{cases}$

$(1)+(2)\Longrightarrow z=0 (4)$

$(4)$ dans $(3)\Longrightarrow x=-y$

$\Rightarrow u=\begin{pmatrix}
-y\\
y\\
0
\end{pmatrix}=y\begin{pmatrix}
-1\\
1  \\
0
\end{pmatrix}\Rightarrow E_0=\langle \begin{pmatrix}
-1\\
1\\
0
\end{pmatrix}\rangle$

$\bullet$ Pour $\lambda=-2$, on a :

$\begin{cases}
2x+2y+2z=-2x(1)\\
-x-y+z=-2y(2)\\
3x+3y+z=-2z(3)
\end{cases}=\begin{cases}
4x+2y+2z=0(1)\\
-x+y+z=0(2)\\
3x+3y+3z=0(3)
\end{cases}=\begin{cases}
2x+y+z=0(1)\\
-x+y+z=0(2)\\
x+y+z=0(3)
\end{cases}$

$(1)-(2)\Rightarrow x=0 (4)$

$(4)$ dans $(3)\Longrightarrow z=-y$

$\Rightarrow u=\begin{pmatrix}
0\\
y\\
-y
\end{pmatrix}=y\begin{pmatrix}
0\\
1\\
-1
\end{pmatrix}\Rightarrow E_{-2}=\langle \begin{pmatrix}
0 \\
1 \\
-1
\end{pmatrix}\rangle$

$\bullet$ Pour $\lambda=4$, on a :

$\begin{cases}
2x+2y+2z=4x(1)\\
-x-y+z=4y(2)\\
3x+3y+z=4z(3)
\end{cases}=\begin{cases}
-2x+2y+2z=0(1)\\
-x-5y+z=0(2)\\
3x+3y-3z=0(3)
\end{cases}=\begin{cases}
-x+y+z=0(1)\\
-x-5y+z=0(2)\\
x+y-z=0(3)
\end{cases}$

$(1)+(3)\Rightarrow y=0 (4)$

$(4)$ dans $(2)\Longrightarrow z=x$

$\Rightarrow u=\begin{pmatrix}
x\\
0\\
x
\end{pmatrix}=x\begin{pmatrix}
1\\
0\\
1
\end{pmatrix}\Rightarrow E_{4}=\langle \begin{pmatrix}
1\\
0\\
1
\end{pmatrix}\rangle$

$P=\begin{pmatrix}
0  & -1 & 1 \\ 
1  & 1  & 0 \\ 
-1 & 0  & 1
\end{pmatrix} $

$det(P)=\begin{vmatrix}
0  & -1 & 1 \\ 
1  & 1  & 0 \\ 
-1 & 0  & 1
\end{vmatrix} = 1+1=2\neq 0$ $P$ est inversible et $P^{-1}=\frac{1}{\det(P)}t_{com(P)}$

$com(P)=\begin{pmatrix}
1  & -1 & 1 \\ 
1  & 1  & 1 \\ 
-1 & 1  & 1
\end{pmatrix}\Rightarrow t_{com(P)}=\begin{pmatrix}
1  & 1 & -1 \\ 
-1  & 1  & 1 \\ 
1 & 1  & 1
\end{pmatrix} \Rightarrow P^{-1}=\frac{1}{2} \begin{pmatrix}
1  & 1 & -1 \\ 
-1  & 1  & 1 \\ 
1 & 1  & 1
\end{pmatrix}$

On vérifie que :

\[PDP^{-1}=\frac{1}{2}\begin{pmatrix}
0  & -1 & 1 \\ 
1  & 1  & 0 \\ 
-1 & 0  & 1
\end{pmatrix}\begin{pmatrix}
-2 & 0& 0 \\ 
0 & 0 & 0 \\ 
0 & 0 & 4
\end{pmatrix}\begin{pmatrix}
1  & 1 & -1 \\ 
-1  & 1  & 1 \\ 
1 & 1  & 1
\end{pmatrix}=\begin{pmatrix}
2 & 2& 2 \\ 
-1 & -1 & 1 \\ 
3 & 3 & 1
\end{pmatrix}=A\]

2.a. Montrons que 
\[
\forall {n\in\mathbb{N}^*}, A^n = PD^nP^{-1}
\]

Raisonnons par récurrence

-Pour $n=1$

D'après 1.b. $PDP^{-1}=A$, donc la propriété est vrai à l'ordre 1.

-Supposons que la propriété est vrai à l'ordre $n$ et vérifions à l'ordre $n+1$

On a : $A^n = PD^nP^{-1}$(*) et  $A=PDP^{-1}$(**).

En multipliant (*) et (**) membres par membres, on obtient :

\begin{align*}
A\times A^n &= PDP^{-1}PD^nP^{-1}\\
A^{n+1} &=PDD^nP^{-1}\\
A^{n+1} &=PD^{n+1}P^{-1}
\end{align*}.
Alors la propriété est vrai à l'ordre $n+1$. On conclut donc que \[
\forall {n\in\mathbb{N}^*}, A^n = PD^nP^{-1}
\]

Déduisons $A^n$.

$A^n = PD^nP^{-1}$ avec $D^n=\begin{pmatrix}
(-2)^n & 0& 0 \\ 
0 & 0 & 0 \\ 
0 & 0 & 4^n
\end{pmatrix}$

$A^n=\frac{1}{2}\begin{pmatrix}
0  & -1 & 1 \\ 
1  & 1  & 0 \\ 
-1 & 0  & 1
\end{pmatrix}\begin{pmatrix}
(-2)^n & 0& 0 \\ 
0 & 0 & 0 \\ 
0 & 0 & 4^n
\end{pmatrix}\begin{pmatrix}
1  & 1 & -1 \\ 
-1  & 1  & 1 \\ 
1 & 1  & 1
\end{pmatrix}$

$A^n=\frac{1}{2}\begin{pmatrix}
4^n & 4^n & 4^n \\ 
(-2)^n & (-2)^n & -(-2)^n \\ 
-(-2)^n+4^n & -(-2)^n+4^n & (-2)^n+4^n
\end{pmatrix}$

2.b. Appliquons ce résultat au  calcul de $A^4$. 

$A^4=\frac{1}{2}\begin{pmatrix}
4^4 & 4^4 & 4^4 \\ 
(-2)^4 & (-2)^4 & -(-2)^4 \\ 
-(-2)^4+4^4 & -(-2)^4+4^4 & (-2)^4+4^4
\end{pmatrix}=\frac{1}{2}\begin{pmatrix}
256 & 256 & 256 \\ 
16 & 16 & -16 \\ 
240 & 240 & 272
\end{pmatrix} =\begin{pmatrix}
128 & 128 & 128 \\ 
8 & 8 & -8 \\ 
120 & 120 & 136
\end{pmatrix}$


\medskip
\textbf{Exercice 2}

Soit l'application : 
\begin{align*}
f:\mathbb{R}^3 &\longrightarrow \mathbb{R}^3\\
(x,y,z) &\longmapsto f(x,y,z)=(-x+y+z, -6x+4y+2z, 3x-y+z)
\end{align*} 

1. Écrivons la matrice de $f$ dans la base canonique de $\mathbb{R}^3$.

Soit $M$ cette matrice,

$M=\begin{pmatrix}
-1 & 1 & 1 \\ 
-6 & 4 & 2 \\ 
3 & -1 & 1
\end{pmatrix}$ 

2. Montrons l'égalité $f\circ f= 2f$.

$f\circ f=f^2=M^2=\begin{pmatrix}
-1 & 1 & 1 \\ 
-6 & 4 & 2 \\ 
3 & -1 & 1
\end{pmatrix}\times\begin{pmatrix}
-1 & 1 & 1 \\ 
-6 & 4 & 2 \\ 
3 & -1 & 1
\end{pmatrix}=\begin{pmatrix}
-2 & 2 & 2 \\ 
-12 & 8 & 4 \\ 
6 & -2 & 2
\end{pmatrix}=2\begin{pmatrix}
-1 & 1 & 1 \\ 
-6 & 4 & 2 \\ 
3 & -1 & 1
\end{pmatrix}=2M=2f$

\medskip
Déduisons que si $v\in Imf$, alors $f(v)=2v$

$v\in Imf\Leftrightarrow \exists u\in \mathbb{R}^3/v=f(u)$

$v=f(u)\Rightarrow f(v)=f(f(u))=f\circ f(u)=f^2(u)=2f(u)=2v$. D'où le résultat.

3. Montrons que les sous-espaces vectoriels $Kerf$ et $Imf$ sont des sous-espaces supplémentaires de $\mathbb{R}^3$

$Kerf$ et $Imf$ sont dites supplémentaire de $\mathbb{R}^3$ ssi $\begin{cases}
Kerf\cap Imf=\{0_{\mathbb{R}^3}\}\\
Kerf+Imf=\mathbb{R}^3
\end{cases}$

Vérifions ces deux conditions.

\medskip
\textbf{Déterminons $Imf$}

D'après 2. si $v\in Imf$, alors $f(v)=2v$

Soit $v=\begin{pmatrix}
x\\
y\\
z
\end{pmatrix}\in Imf$, alors $f(v)=2v$

$f(v)=2v \Rightarrow \begin{cases}
-x+y+z=2x\\
-6x+4y+2z=2y\\
3x-y+z=2z
\end{cases}\Rightarrow \begin{cases}
-3x+y+z=0\\
-6x+2y+2z=0\\
3x-y-z=0
\end{cases}\Rightarrow \begin{cases}
-3x+y+z=0\\
-3x+y+z=0\\
-3x+y+z=0
\end{cases}\Rightarrow z=3x-y$

$\Rightarrow v=\begin{pmatrix}
x\\
y\\
z
\end{pmatrix}=\begin{pmatrix}
x\\
y\\
3x-y
\end{pmatrix}=\begin{pmatrix}
x\\
0\\
3x
\end{pmatrix}+\begin{pmatrix}
0\\
y\\
-y
\end{pmatrix}=x\begin{pmatrix}
1\\
0\\
3
\end{pmatrix}+y\begin{pmatrix}
0\\
1\\
-1
\end{pmatrix}$

D'où $Imf=\langle \begin{pmatrix}
1\\
0\\
3
\end{pmatrix},\begin{pmatrix}
0 \\
1 \\
-1
\end{pmatrix}\rangle$ et $dimImf=2$

\medskip
\textbf{Déterminons $Kerf$}

$Kerf=\{u=\begin{pmatrix}
x\\
y\\
z
\end{pmatrix}\in \mathbb{R}^3/f(u)=0_{\mathbb{R}^3}\}$

$f(u)=0_{\mathbb{R}^3}\Rightarrow \begin{cases}
-x+y+z=0\\
-6x+4y+2z=0\\
3x-y+z=0
\end{cases}\Rightarrow \begin{cases}
-x+y+z=0(1)\\
-3x+2y+z=0(2)\\
3x-y+z=0(3)
\end{cases}$

$(2)+(3)\Rightarrow y=-2z$ (4)

(4) dans (1) $\Rightarrow x=-z$

$\Rightarrow u=\begin{pmatrix}
x\\
y\\
z
\end{pmatrix}=\begin{pmatrix}
-z\\
-2z\\
z
\end{pmatrix}=z\begin{pmatrix}
-1 \\
-2 \\
1
\end{pmatrix}$

D'où $Kerf=\langle \begin{pmatrix}
-1\\
-2\\
1
\end{pmatrix}\rangle$ et $dimKerf=1$.

\medskip
On a : $dimImf+dimKerf=2+1=3=dim\mathbb{R}^3$ et $Imf+Kerf=\mathbb{R}^3$, donc la deuxième condition est vérifiée.

\medskip
Soit $N=\begin{pmatrix}
1 & 0  & -1 \\ 
0 & 1  & -2 \\ 
3 & -1 & 1 
\end{pmatrix}$ la matrice formée par les vecteurs de $Imf$ et de $Kerf$

On a $\det(N)=\begin{vmatrix}
1 & 0  & -1 \\ 
0 & 1  & -2 \\ 
3 & -1 & 1 
\end{vmatrix}=(1-2)-(0-3)=-1+3=2\neq 0$ alors la famille formée par les trois vecteurs de $Imf$ et $Kerf$ est libre.

D'où $Imf\cap Kerf=\{0_{\mathbb{R}^3}\}$. La deuxième condition est donc vérifiée.

\medskip
Ces deux conditions étant vérifiées, on conclut donc que les sous-espaces vectoriels $Kerf$ et $Imf$ sont des sous-espaces supplémentaires de $\mathbb{R}^3$


4. Trouvons une base $(e_1,e_2,e_3)$ de $\mathbb{R}^3$ dont le premier vecteur appartient à $Kerf$ et les derniers à $Imf$. 

$e_1\in Kerf\Rightarrow f(e_1)=0_{\mathbb{R}^3}\Longrightarrow e_1=\begin{pmatrix}
-1\\
-2\\
1
\end{pmatrix}$ (Calculé dans la question 3)

$e_2,e_3\in Imf \Longrightarrow e_2=\begin{pmatrix}
1\\
0\\
3
\end{pmatrix}$ et $e_3=\begin{pmatrix}
0 \\
1 \\
-1
\end{pmatrix}$

D'où $(e_1,e_2,e_3)=(\begin{pmatrix}
-1\\
-2\\
1
\end{pmatrix},\begin{pmatrix}
1\\
0\\
3
\end{pmatrix},\begin{pmatrix}
0 \\
1 \\
-1
\end{pmatrix})$

\medskip
La matrice de $f$ dans la base $(e_1,e_2,e_3)$. Soit 	$A$ cette matrice.

$f(e_1)=(1-2+1,6-8+2,-3+2+1)=(0,0,0)$

$f(e_2)=(-1+0+3,-6+0+6,3-0+3)=(2,0,6)=2e_2$

$f(e_3)=(-0+1-1,-0+4-2,0-1-1)=(0,2,-2)=2e_3$

$A=\begin{pmatrix}
0 & 0 & 0 \\ 
0 & 2 & 0 \\ 
0 & 0 & 2
\end{pmatrix} $


5. Trouvons une équation de $Imf$.

D'après 2. si $v\in Imf$, alors $f(v)=2v$

Soit $v\in Imf\Longrightarrow f(v)=2v\Longrightarrow Mv=2v\Longrightarrow Mv-2v=0\Longrightarrow (M-2I_3)v=0$

Ainsi $Imf=\{v\in \mathbb{R}^3/(M-2I_3)v=0\}$

\medskip
\textbf{Exercice 3}

On pose pour $(x,y)\in\Omega=\mathbb{R}\times]0,+\infty[$,  $f(x,y)=y(x^2 + ln^2{(y)})$.

1.a. Calculons $\frac{\partial{f}}{\partial{x}}(x,y)$ 

$\frac{\partial{f}}{\partial{x}}(x,y)=2xy$

\medskip
trouvons les $(x,y)\in\Omega$ tels que $\frac{\partial{f}}{\partial{x}}(x,y) = 0$.

$\frac{\partial{f}}{\partial{x}}(x,y) = 0\Rightarrow 2xy=0\Rightarrow x=0$ ou $y=0$

Comme $x\in\mathbb{R}$ et $y\in ]0,+\infty[$ alors $y>0$ et $x=0$.

Donc l'ensemble recherché est $\{(0,y), \forall y\in ]0,+\infty[\}$


1.b. Calculons $\frac{\partial{f}}{\partial{y}}(x,y)$éduire les points critiques de $f$.

$\frac{\partial{f}}{\partial{y}}(x,y)=x^2 + ln^2{(y)}+2ln(y)$

\medskip
Déduisons les points critiques de $f$.

Les points critiques sont les points qui annulent les dérivées premières.

$\begin{cases}
\frac{\partial{f}}{\partial{x}}(x,y) = 0 (1)\\
\frac{\partial{f}}{\partial{y}}(x,y) = 0 (2)
\end{cases}$

De (1) $\frac{\partial{f}}{\partial{x}}(x,y) = 0\Longrightarrow x=0$ d'après la question 1.a)

De (2) $\frac{\partial{f}}{\partial{y}}(x,y) = 0$ et $x=0$ $\Rightarrow ln^2{(y)}+2ln(y)=0\Rightarrow ln(y)(ln(y)+2)=0\Rightarrow ln(y)=0$ ou $ln(y)+2=0$

$ln(y)=0\Rightarrow y=1$

$ln(y)+2=0\Rightarrow ln(y)=-2\Rightarrow y=e^{-2}$

Donc les points critiques sont : $A=(0,1)$ et $B=(0,e^{-2})$



2.a. Déduisons les extrema relatifs de $f$.

Calculons la matrice Hessienne de $f$, $Hess_f(x,y)$

$Hess_f(x,y)=\begin{pmatrix}
\frac{\partial^2{f}}{\partial{x^2}}(x,y) & \frac{\partial^2{f}}{\partial{x}\partial{y}}(x,y) \\ 
\frac{\partial^2{f}}{\partial{y}\partial{x}}(x,y) & \frac{\partial^2{f}}{\partial{y^2}}(x,y)
\end{pmatrix}$

$\frac{\partial^2{f}}{\partial{x^2}}(x,y)=2y$, $\frac{\partial^2{f}}{\partial{y^2}}(x,y)=\frac{2}{y}(ln(y)+1)$, $\frac{\partial^2{f}}{\partial{x}\partial{y}}(x,y)=\frac{\partial^2{f}}{\partial{x}\partial{y}}(x,y)=2x$

$\Longrightarrow Hess_f(x,y)=\begin{pmatrix}
2y & 2x \\ 
2x & \frac{2}{y}(ln(y)+1)
\end{pmatrix}$

$\bullet$ Nature de $A=(0,1)$

$Hess_f(A)=\begin{pmatrix}
2 & 0 \\ 
0 & 2
\end{pmatrix} \Longrightarrow det(Hess_f(A))=\begin{vmatrix}
2 & 0 \\ 
0 & 2
\end{vmatrix}=4-0=4>0$ et $\frac{\partial^2{f}}{\partial{x^2}}(A)=2>0$ alors $A$ est un minimum local.

$\bullet$ Nature de $B=(0,e^{-2})$

$Hess_f(B)=\begin{pmatrix}
2e^{-2} & 0 \\ 
0 & -2e^2
\end{pmatrix} \Longrightarrow det(Hess_f(B))=\begin{vmatrix}
2e^{-2} & 0 \\ 
0 & -2e^{2}
\end{vmatrix}=-4-0=-4<0$ alors $B$ est un point selle.


2.b. Montrons que si $(x,y)\in\Omega$ et $(x,y)\neq(0,1)$, on a $f(x,y)>0$.

$\forall (x,y)\in\Omega$ et $(x,y)\neq(0,1)$, on a :

$y>0$, $x^2>0$ et $ln^2(y)>0$ $\Rightarrow y(x^2+ln^2(y))>0\Rightarrow f(x,y)>0$

\medskip
On a : $f(0,1)=1(0^2+ln^2(1))=0$

Comme $\begin{cases}
\forall (x,y)\in\Omega\hspace{0.7cm} et\hspace{0.7cm} (x,y)\neq(0,1), f(x,y)>0\\
et\\
f(0,1)=0
\end{cases}$ alors on en déduire que (0,1) est un minimum global.

\medskip\textbf{Exercice 4}

a. Soit $\alpha\in \mathbb{R}$. Résolvons l'équation $(E_1):y^{\prime\prime} -2y^{\prime} +(1-\alpha^2)y=e^x$.

Cherchons la solution de l'équation homogène $(E_0) : y^{\prime\prime} -2y^{\prime} +(1-\alpha^2)y=0$

L'équation caractéristiques est :

$r^2-2r^2+1-\alpha^2=0$

$\Delta = 4-4(1-\alpha^2)=4\alpha^2=(2\alpha)^2$

$r_1=\frac{2-2\alpha}{2}=1-\alpha$ et $r_2=\frac{2+2\alpha}{2}=1+\alpha$

Nous avons deux cas :

$\bullet$ Si $\alpha=0$, on a : $r_1=r_2=1$

Ainsi, $y_0(x)=(Ax+B)e^x$

\textbf{Solution particulière :}

Le second membre $\varphi (x)=e^x$ est de la forme $e^{mx}$, avec  $m=1=r_1=r_2$, il existe donc une solution particulière sous la forme $y_p(x)=Cx^2e^x$ avec $C$ une constante réelle à déterminer. On a:

$y_p(x)=Cx^2e^x$

$y'_p(x)=2Cxe^x+Cx^2e^x=(2Cx+Cx^2)e^x$

$y''_p(x)=(2C+4Cx+Cx^2)e^x$

En reportant dans l'équation différentielle, on obtient :

$(2C+4Cx+Cx^2-2(2Cx+Cx^2)+Cx^2)e^x=e^x$

$2C=1\Rightarrow C=\frac{1}{2}$

Donc la solution particulière est $y_p(x)=\frac{1}{2}x^2e^x$

Ainsi la solution générale de l'équation pour $\alpha=0$ est  

$y(x)=y_0(x)+y_p(x)=(Ax+B)e^x+ \frac{1}{2}x^2e^x=(\frac{1}{2}x^2+Ax+B)e^x$ avec $A,B\in \mathbb{R}$

\medskip
$\bullet$ Pour $\alpha\neq 0$

L'équation caractéristique a deux racines distinctes $r_1=1-\alpha$ et $r_2=1+\alpha$

Donc la solution de l'équation homogène est donnée par :

$y_0(x)=C_1e^{(1-a)x}+C_2e^{(1+a)x}$

\medskip
Cherchons une solution particulière

Le second membre $\varphi (x)=e^x$ est de la forme $e^{mx}$ ou $m=1$ n'est pas une racine de l'équation caractéristique.

Donc il existe une solution particulière de la forme $y_p(x)=Ke^x$ avec $K$ une constante réelle à déterminer.

On a :

$y_p(x)=Ke^x$

$y'_p(x)=Ke^x$

$y''_p(x)=Ke^x$

En remplaçant dans l'équation, on obtient :

$K-2K+(1-\alpha^2)K)e^x=e^x$

$-K\alpha^2=1\Longrightarrow K=-\frac{1}{\alpha^2}$

Donc la solution particulière est $y_p(x)=-\frac{1}{\alpha^2}e^x$

Ainsi la solution générale de l'équation différentielle pour $\alpha\neq 0$ est :

$y(x)=C_1e^{(1-a)x}+C_2e^{(1+a)x}-\frac{1}{\alpha^2}e^x$.

b. Déduisons la solution de l'équation différentielle $(E_2):y^{\prime\prime} -2y^{\prime}+y=e^x$ satisfaisant au conditions initiales $y(0)=1$ et $y^\prime(0)=3$.


La solution générale de l'équation est :

$(\frac{1}{2}x^2+Ax+B)e^x$ d'après la première question.

$y(0)=1\Rightarrow B=1$

$y'(x)=(x+A+\frac{1}{2}x^2+Ax+B)e^x=(\frac{1}{2}x^2+(A+1)x+A+B)e^x$

$y'(0)=3\Longrightarrow A+B=3\Longrightarrow A=2$

Donc la solution recherchée est $y(x)=(\frac{1}{2}x^2+2x+1)e^x$

\medskip

\end{spacing}
\end{document}
